\section{Proof of the local counting and fixing lemmas}
\label{app:local}

Throughout this appendix the active graph is acyclic and fixed items have
been absorbed into propagated intervals. All laws are uniform on the finite
set of admissible weak assignments. The random variables used in the proof
are not sampled by the adversary.

\subsection{Compression and insertion}

\begin{proof}[Proof of Lemma~\ref{lem:insertion}]
Replace stored intervals by effective intervals, which leaves the family
unchanged. Delete $x$, keeping those intervals and all implied inequalities
between remaining items. Let $N^-$ count the resulting assignments and $N$
the original ones.

Every deleted assignment admits an insertion of $x$: its predecessors lie
before its successors, predecessor values are at most $r$, and successor
values are at least $\ell$. The allowed positions are a nonempty interval.
There are at most $W$ of them, so
\begin{equation}
                              N\leq WN^-.
\label{eq:delete-upper}
\end{equation}

Transform a deleted assignment as follows. For each proper predecessor of
$x$, leave values below $\ell$ unchanged and replace
$\ell+i$ by $\ell+\lfloor i/4\rfloor$ for $i\geq0$.
For each proper successor, leave values above $r$ unchanged and replace
$r-i$ by $r-\lfloor i/4\rfloor$.
Other coordinates are unchanged.

These transformations preserve all intervals: predecessor lower bounds are
at most $\ell$, and successor upper bounds at least $r$.
They are nondecreasing within each class. No inequality enters the
predecessor set from outside it, while lowering a predecessor helps an
inequality leaving that set. The symmetric statement holds for successors.
Thus every transformed assignment is admissible in the deleted system.

There are at most $4^c$ preimages of any transformed assignment.
All transformed predecessors lie at or below $\ell+t_0$, and successors
at or above $r-t_0$. Each position in $I_x$ can therefore receive $x$.
It follows that
\begin{equation}
                    N\geq4^{-c}|I_x|N^-\geq \frac{WN^-}{2\cdot4^c}.
\label{eq:delete-lower}
\end{equation}
For a fixed $t\in I_x$, at least $4^{-c}N^-$ assignments have $X_x=t$.
Divide by \eqref{eq:delete-upper} to prove \eqref{eq:central}.
An interval of $m$ positions allows at most $m$ insertions for each deleted
assignment; divide $mN^-$ by \eqref{eq:delete-lower} to prove
\eqref{eq:concentration}. With a different coordinate $Y$ fixed in the
deleted assignment, at most $2z+1$ integer positions satisfy
$|X_x-Y|\leq z$. This proves \eqref{eq:near-tie}, including dependent
coordinates and fixed real thresholds.
\end{proof}

\subsection{Conditional refresh and a one-sweep bound}

Take a topological order of the active graph. One sweep visits coordinates
in that order and replaces each by a uniform allowed integer conditional
on all other coordinates. Its allowed set is an interval $[a,b]$, and the
replacement can be written, for a fresh uniform $U\in[0,1)$,
\begin{equation}
                   \lfloor a+U(b-a+1)\rfloor
                   =\lfloor(1-U)a+Ub+U\rfloor.
\label{eq:refresh}
\end{equation}
Each conditional update preserves the uniform assignment law. Thus starting
in that law and performing any number of sweeps leaves it unchanged.

Compare full and truncated systems using the same fresh $U$ at every common
coordinate, and restrict the full topological order for the local sweep.
Let $\Delta_x$ be the absolute discrepancy of their coordinate values.
At a coordinate whose incident edges are present in both systems,
\begin{equation}
\Delta_x^{\mathrm{new}}
\leq1+(1-U)\max_{y\to x}\Delta_y^{\mathrm{new}}
          +U\max_{x\to z}\Delta_z^{\mathrm{old}}.
\label{eq:one-update}
\end{equation}
Empty maxima are zero. The lower conditional endpoint is the maximum
of its fixed interval bound and predecessor values; the upper endpoint
is the minimum of its fixed interval bound and successor values.
Equation~\eqref{eq:refresh} gives \eqref{eq:one-update}, with one unit
for integer rounding.

Assume height at most $h$. Expand the predecessor terms backward through
the sweep. Within at most $h$ updates each predecessor-only path reaches
a common fixed lower bound. Let $A_x$ be the minimum of the fresh
$1-U$ values at $x$ and all its active predecessors.
The coefficient reaching such a common bound is at least $A_x^h$.
At most $h$ rounding errors accumulate, and any old discrepancy left in
the expansion lies within undirected graph distance $h$ of $x$.
Consequently
\begin{equation}
 \Delta_x^{\mathrm{new}}\leq
 h+(1-A_x^h)\max_{y\in B_h(x)}\Delta_y^{\mathrm{old}}.
\label{eq:sweep}
\end{equation}
The minimum defining $A_x$ covers every possible maximizing predecessor,
so this inequality does not select a preferred path in advance.

\subsection{Moment contraction in a shrinking interior}

For maximum degree $d$, put
\begin{equation}
 M=(d+1)^{h+1},\qquad a_0=\frac1{8M^2},\qquad
 s=\left\lceil\frac{8M}{a_0^h}\right\rceil .
\label{eq:moment-constants}
\end{equation}
Both $B_h(x)$ and the set of fresh numbers defining $A_x$ have at most $M$
members. A union bound gives $\Prob(A_x<a_0)\leq Ma_0$.
On the complementary event $(1-A_x^h)^s\leq e^{-sa_0^h}$, and hence
\begin{equation}
            M\,\E(1-A_x^h)^s
            \leq M(Ma_0+e^{-sa_0^h})\leq\frac12.
\label{eq:moment-contraction}
\end{equation}

Retain the radius-$ht$ neighborhoods of the requested coordinates.
Start each system in its own uniform law, with any initial coupling, and
run $t$ coupled sweeps. At sweep $j$ restrict attention to
$S_j$, the union of the radius-$h(t-j)$ neighborhoods of the requested
coordinates. For $x\in S_j$, $B_h(x)\subseteq S_{j-1}$, so all edges
needed for \eqref{eq:sweep} are present.

Writing $z_+=\max\{z,0\}$, \eqref{eq:sweep} implies
\[
 (\Delta_x^{\mathrm{new}}-hj)_+
 \leq(1-A_x^h)
       \max_{y\in B_h(x)}(\Delta_y^{\mathrm{old}}-h(j-1))_+.
\]
The fresh numbers are independent of the old discrepancies.
Raise to the power $s$, bound a maximum by the sum of the $s$th powers,
and apply \eqref{eq:moment-contraction}. The largest moment in the
shrinking interior falls by a factor at least two per sweep.
Initially, both values at a common coordinate lie in its same globally
propagated interval. If $W_{\max,x}$ is the largest width in $x$'s full
component, induction gives at a requested coordinate
\begin{equation}
                \E(\Delta_x-ht)_+^s\leq W_{\max,x}^s2^{-t}.
\label{eq:moment-final}
\end{equation}
No assumption about the truncated boundary law has been made.

\subsection{Explicit radius and minimum width}
\label{app:radius}

There are at most $2M$ other coordinates comparable with one coordinate.
For a desired error $\varepsilon$, set
\begin{equation}
             C_{\mathrm{ins}}=2\cdot4^{2M},\qquad
             u=\frac{\varepsilon}{16C_{\mathrm{ins}}}.
\label{eq:prob-constants}
\end{equation}
Choose an integer $t$ and integer $W_0$ satisfying
\begin{align}
t&\geq s\logtwo(2R/u)+\logtwo(8/\varepsilon),\label{eq:t}\\
W_0&\geq\max\{32,\;2ht/u,\;4C_{\mathrm{ins}}/\varepsilon\}.
\label{eq:w0}
\end{align}
Take $\rho=ht$. These choices are computable with integer upper roundings.

For a nonisolated requested coordinate, $W_{\max,x}\leq RW_x$ and
$ht\leq uW_x/2$. Markov's inequality and \eqref{eq:moment-final} give
\begin{equation}
       \Prob(\Delta_x>uW_x)
       \leq(2R/u)^s2^{-t}\leq\varepsilon/8.
\label{eq:coordinate-error}
\end{equation}
Outside these exceptional events, disagreement between the two weak
comparison indicators requires
$|X_x-X_y|\leq u(W_x+W_y)$. Applying \eqref{eq:near-tie} at the
larger-width endpoint, of width $W$, bounds this probability by
\[
                  4C_{\mathrm{ins}}u+C_{\mathrm{ins}}/W
                  \leq\varepsilon/2.
\]
The two exceptional probabilities leave a total less than $\varepsilon$.
A fixed threshold uses the same argument with one requested coordinate.
An isolated coordinate is exactly uniform and independent of all others:
couple its values identically. If both coordinates are isolated the local
answer is exact; if one is nonisolated the larger-width endpoint in
\eqref{eq:near-tie} still has width at least $W_0$.
Both experiments preserve their uniform laws, so these disagreement bounds
prove Lemma~\ref{lem:local}.

\subsection{Log-concavity by an injection}
\label{app:logconcavity}

Let $p_i=\Prob(X_x=i)$. We prove
\begin{equation}
                              p_{i-1}p_{i+1}\leq p_i^2.
\label{eq:logconcavity}
\end{equation}
Take admissible assignments $f,g$ with $f(x)=i-1$, $g(x)=i+1$,
and write $h=g-1$ coordinatewise. Both $f$ and $h$ respect all weak
edges, although $h$ may violate interval bounds.
On the items form an undirected graph containing the edge of an inequality
$y\to z$ exactly when
\begin{equation}
                  \max\{f(y),h(y)\}>\min\{f(z),h(z)\}.
\label{eq:swap-graph}
\end{equation}
Across such an edge $h-f$ is nonzero and has the same sign at both ends.
For example, if $h(y)>f(y)$ and $h(z)\leq f(z)$, then
$h(y)\leq h(z)$ contradicts \eqref{eq:swap-graph}. The other sign change
is symmetric.

On the component containing $x$ we therefore have $h>f$.
Swap $f,h$ on that component, then add one to the second assignment.
Both output assignments take value $i$ at $x$.
For an item $y$ in the component,
$\ell_y\leq f(y)<h(y)\leq r_y-1$, so the swapped values, with the
indicated addition, obey its original interval.
An edge crossing the component boundary fails
\eqref{eq:swap-graph}, so either choice of its two values preserves
the inequality. Other edges are unchanged or swapped together.
The outputs are therefore admissible.

The map is injective. Subtract one from the second output: the unordered
coordinate pairs are precisely the original $\{f(y),h(y)\}$.
They determine \eqref{eq:swap-graph} and its component at $x$.
Swap back and add one to recover $g$. Counting the input and output
pairs proves \eqref{eq:logconcavity}.

By Lemma~\ref{lem:support}, positive probabilities occupy an interval.
The ratios $p_{i+1}/p_i$ decrease there. The cumulative sums $F_i$ are
log-concave too: for $k\leq i-1$,
$p_kp_{i+1}\leq p_{k+1}p_i$. Summing gives
$F_{i-1}p_{i+1}\leq F_ip_i$, equivalently
$F_{i-1}F_{i+1}\leq F_i^2$. Reflection gives the same property for
inclusive tails. Zero terms at the endpoints cause no difficulty.

\subsection{Approximate quartiles give a reliable interval}
\label{app:quartiles}

\begin{proof}[Proof of Lemma~\ref{lem:quartiles}]
Since $\|\widehat F-F\|_\infty\leq1/16$ and $a,b$ are first crossings,
\[
\begin{array}{ll}
 F(a)\geq3/16, & F(a-1)<5/16,\\
 F(b)\geq11/16,& F(b-1)<13/16.
\end{array}
\]
Put $m=b-a+1$. If $F(a-1)<1/8$, then $p_a>1/16$.
Otherwise log-concavity of $F$ implies
\[
        \frac{11}{5}\leq\frac{F(b)}{F(a-1)}
                   \leq\left(\frac{F(a)}{F(a-1)}\right)^m,
\]
and therefore
\[
 p_a\geq\frac18\bigl((11/5)^{1/m}-1\bigr)
      \geq\frac{\ln(11/5)}{8m}>\frac1{16m}.
\]
The reflected argument using inclusive tails gives $p_b\geq1/(16m)$.
For $m>1$, log-concavity of $p$ gives, for $0\leq r<m$,
\[
 p_{a+r}\geq p_a^{\,1-r/(m-1)}p_b^{\,r/(m-1)}
           \geq\min\{p_a,p_b\}.
\]
The case $m=1$ already follows from the endpoint bound.
The local coordinate law stays within the original effective interval,
whose full support is attained globally, so its crossings lie there.
\end{proof}
